\scanpage{55}
\uline{Question 1} Supposons que \deleted{$G$ ait un} $T$ qui \deleted{est}\inserted{soit} $\simeq T_1$, alors ce $T$ est-il contenu dans un Borel, i.e. est-ce qu'il existe un isom \uline{intérieur} $T_1-\unsure{T_\rho|T}$ ??
J'en doute.

\uline{Question 2} Supposons que \deleted{$T_1$ soit isom. à} $G$ ait un \deleted{\unsure{groupe}} tore maximal $T$ contenu dans un Borel. S'ensuit-il que $G\simeq G_1$ ? \deleted{\unsure{Je ne sais}} \deleted{\unsure{si on a} $\rho$ \unsure{trivial}} C'est le cas si $\operatorname{Pic}(S)=0$ \inserted{(p. ex. $S$ local)} et si $\rho$ \deleted{\unsure{est}} est trivial. Est-ce vrai sur un corps $k$ ($\rho$ non trivial) ??

\margin{N.B. on utilise que le normalisateur de $(B_0,T_0)$ dans $\Gamma_0$ est $T'_0$ \deleted{\unsure{mais} $\unsure{T_0}$} \deleted{\illegible{3}} \unsure{alors} \unsure{dans le cas} \unsure{d'un groupe} \unsure{on a} \deleted{\unsure{un groupe}} \unsure{``} $\unsure{W}$ \unsure{''} \deleted{\illegible{2}} \unsure{extérieure}.}

\deleted{Par exemple $k=\mathbb R$, groupe unitaire.} C'est encore vrai si $k$ est un \uline{corps de Tsen}
[$H^1(k,T)=0$ pour \uline{$\forall$} tore sur $k$]. Que se passe-t-il pour $k=\mathbb R$, \deleted{\unsure{lorsque}} $G_1,G$ groupes \inserted{projectifs} unitaires spécial à deux variables ???

\margin{Dans ce cas on aura $\xi=e$}

\deleted{\unsure{\S}}
\uline{Question 3} Supposons qu'il existe un sous-groupe de Borel $B$ de $G$ \inserted{\unsure{localement conjugué} \unsure{au} $G_1:B_1$}, Alors peut-on réduire le groupe structural \deleted{à} de $G_1$ à $T_1$, et $=$ est-il vrai que \deleted{\unsure{l'on}} l'on peut trouver un sous-groupe
\margin{54}
